%! AP Statistics — Unit 5, Lesson 5.2 — Warm-Up (Answer Key)
\documentclass[10pt]{article}
\usepackage{apstats-article}
\usepackage{apstats-key}   % pulls in -boxes; do NOT also load -boxes
\newcommand{\TallMath}[1]{$\displaystyle #1\rule[-1.4em]{0pt}{3.2em}$}

\begin{document}

\pageheader{Unit 5, Lesson 5.2}{Warm-Up --- Answer Key}
\namedateperiod

\vspace{0.15in}

\begin{notesbox}{Spiral Review: Normal Distributions --- Key}
\textbf{1.} Adult male heights are approximately normally distributed with mean
$\mu = 69.2$ inches and standard deviation $\sigma = 3.6$ inches.

\begin{enumerate}[label=\alph*.]
  \item Compute the z-score for a height of 75 inches.
        \ans{$z = (75 - 69.2)/3.6 \approx 1.61$}
  \item Using the 68--95--99.7 rule, what percentage of adult males are between
        62 and 76.4 inches tall?
        \ans{62 = $\mu - 2\sigma$ and $76.4 = \mu + 2\sigma$, so approximately 95\%.}
  \item Would a height of 80 inches be unusual? Justify using a z-score.
        \ansline{$z = (80 - 69.2)/3.6 \approx 3.0$. Yes, this is 3 standard deviations above the mean --- only about 0.15\% of males are this tall, so it would be very unusual.}
\end{enumerate}

\vspace{0.1in}
\textbf{2.} \textit{(Bell-shaped curve centered at 100; students should mark 85, 70, 55 to the left and 115, 130, 145 to the right.)}

\vspace{1.4in}
\end{notesbox}

\end{document}
